\chapter{Supplementary Materials}

This appendix provides supplementary mathematical derivations, additional results tables, and the declaration of AI usage.

\section{Derivation of the IRLS Algorithm}

The Newton-Raphson update for logistic regression is:

\[
\boldsymbol{\beta}^{(t+1)} = \boldsymbol{\beta}^{(t)} - \left(\frac{\partial^2 \ell}{\partial \boldsymbol{\beta} \partial \boldsymbol{\beta}^\top}\right)^{-1} \frac{\partial \ell}{\partial \boldsymbol{\beta}}
\]

Substituting the score and Hessian:

\[
\boldsymbol{\beta}^{(t+1)} = \boldsymbol{\beta}^{(t)} + (\mathbf{X}^\top \mathbf{W}^{(t)} \mathbf{X})^{-1} \mathbf{X}^\top (\mathbf{y} - \mathbf{p}^{(t)})
\]

Let $\mathbf{z}^{(t)} = \mathbf{X}\boldsymbol{\beta}^{(t)} + (\mathbf{W}^{(t)})^{-1}(\mathbf{y} - \mathbf{p}^{(t)})$. Then:

\[
\boldsymbol{\beta}^{(t+1)} = (\mathbf{X}^\top \mathbf{W}^{(t)} \mathbf{X})^{-1} \mathbf{X}^\top \mathbf{W}^{(t)} \mathbf{z}^{(t)}
\]

This is the IRLS update.

\section{Derivation of the ICC for Logistic GLMM}

For a logistic GLMM, the level-1 residual variance is fixed at $\pi^2/3 \approx 3.28987$ (variance of the standard logistic distribution). The between-cluster variance is $\sigma_u^2$.

The intra-cluster correlation is the proportion of total variance explained by between-cluster differences:

\[
\operatorname{ICC} = \frac{\sigma_u^2}{\sigma_u^2 + \pi^2/3}
\]

This follows from decomposing the total variance into between-cluster and within-cluster components.

\section{Derivation of the E-value}

For a protective effect with $RR < 1$, define $RR_{\text{inv}} = 1/RR > 1$. The E-value is the solution to:

\[
RR_{\text{inv}} = \text{E-value} + \sqrt{\text{E-value} \times (\text{E-value} - 1)}
\]

Let $x = \text{E-value}$. Then:

\[
RR_{\text{inv}} = x + \sqrt{x(x-1)}
\]

Let $y = \sqrt{x(x-1)}$. Then $y^2 = x^2 - x$, and $RR_{\text{inv}} - x = y$. Squaring both sides:

\[
(RR_{\text{inv}} - x)^2 = x^2 - x
\]

\[
RR_{\text{inv}}^2 - 2RR_{\text{inv}} x + x^2 = x^2 - x
\]

\[
RR_{\text{inv}}^2 - 2RR_{\text{inv}} x = -x
\]

\[
RR_{\text{inv}}^2 = 2RR_{\text{inv}} x - x = x(2RR_{\text{inv}} - 1)
\]

\[
x = \frac{RR_{\text{inv}}^2}{2RR_{\text{inv}} - 1} = RR_{\text{inv}} + \sqrt{RR_{\text{inv}}(RR_{\text{inv}} - 1)}
\]

This is the E-value formula.

\section{Additional Results Tables}

\subsection{PSM Results by Caliper}

Table~\ref{tab:appendix_psm} shows PSM results for all calipers tested.

\begin{table}[htbp]
	\centering
	\caption{PSM Results by Caliper}
	\label{tab:appendix_psm}
	\begin{tabular}{@{}lcccc@{}}
		\toprule
		\textbf{Survey} & \textbf{Caliper} & \textbf{OR} & \textbf{95\% CI} & \textbf{Matched N} \\
		\midrule
		KMIS 2015 & 0.2 & 0.533 & (0.282-1.010) & 1,742 \\
		KMIS 2015 & 0.1 & 0.456 & (0.233-0.891) & 1,592 \\
		KMIS 2015 & 0.05 & 0.392 & (0.194-0.791) & 1,508 \\
		KMIS 2020 & 0.2 & 0.442 & (0.242-0.810) & 2,350 \\
		KMIS 2020 & 0.1 & 0.465 & (0.242-0.894) & 2,230 \\
		KMIS 2020 & 0.05 & 0.466 & (0.242-0.896) & 2,178 \\
		\bottomrule
	\end{tabular}
\end{table}

\subsection{IPTW Weight Truncation}

Table~\ref{tab:appendix_iptw} shows IPTW truncation details.

\begin{table}[htbp]
	\centering
	\caption{IPTW Weight Truncation}
	\label{tab:appendix_iptw}
	\begin{tabular}{@{}lccc@{}}
		\toprule
		\textbf{Survey} & \textbf{Original Max Weight} & \textbf{Truncated Obs} & \textbf{OR after Truncation} \\
		\midrule
		KMIS 2015 & 9.74 & 13 & 0.482 \\
		KMIS 2020 & 3.89 & 0 & 0.511 \\
		\bottomrule
	\end{tabular}
\end{table}

\section{Summary of All Generated Figures}

Table~\ref{tab:appendix_figures} provides a complete list of all figures generated in the analysis.

\begin{table}[htbp]
	\centering
	\caption{Complete List of Generated Figures}
	\label{tab:appendix_figures}
	\begin{tabular}{@{}ll@{}}
		\toprule
		\textbf{Figure} & \textbf{Description} \\
		\midrule
		dag\_causal\_diagram\_base.png & Base R DAG visualization \\
		dag\_causal\_diagram\_ggraph.png & ggraph DAG visualization \\
		glmm\_caterpillar\_*.png & GLMM random effects caterpillar plots \\
		psm\_love\_plot\_*.png & Love plots for PSM balance assessment \\
		iptw\_weights\_*.png & IPTW weight distributions \\
		dr\_potential\_outcomes\_*.png & Doubly robust potential outcomes plots \\
		ps\_overlap\_*.png & Propensity score overlap plots \\
		week4\_odds\_ratios\_comparison.png & Forest plot comparing traditional methods \\
		week5\_dml\_forest\_ate.png & Forest plot of DML ATE estimates \\
		week5\_evalue\_plot.png & E-value sensitivity analysis plot \\
		week6\_heterogeneity\_forest.png & Heterogeneity forest plot \\
		week6\_robustness\_comparison.png & Robustness comparison across learners \\
		malaria\_by\_wealth.png & Malaria prevalence by wealth quintile \\
		itn\_by\_wealth.png & ITN use by wealth quintile \\
		malaria\_by\_residence.png & Malaria prevalence by urban/rural \\
		malaria\_by\_age.png & Malaria prevalence by age group \\
		itn\_by\_age.png & ITN use by age group \\
		malaria\_map\_2015.png & Geographic map of malaria prevalence (2015) \\
		malaria\_map\_2020.png & Geographic map of malaria prevalence (2020) \\
		itn\_map\_2015.png & Geographic map of ITN use (2015) \\
		itn\_map\_2020.png & Geographic map of ITN use (2020) \\
		correlation\_matrix.png & Correlation matrix of key variables \\
		\bottomrule
	\end{tabular}
\end{table}

\section{Computational Environment}

The analysis was conducted using R version 4.2.2. Key packages used:

\begin{verbatim}
	- dplyr (1.1.0)      : Data manipulation
	- tidyr (1.3.0)      : Data cleaning
	- ggplot2 (3.4.1)    : Visualizations
	- survey (4.2)       : Survey-weighted analysis
	- lme4 (1.1-31)      : GLMM estimation
	- MatchIt (4.5.1)    : Propensity score matching
	- cobalt (4.4.1)     : Balance assessment
	- DoubleML (0.5.0)   : Double Machine Learning
	- mlr3 (0.15.0)      : Machine learning framework
	- ranger (0.14.1)    : Random forest
	- xgboost (1.7.3.1)  : Gradient boosting
\end{verbatim}

All code used to generate the results in this thesis is available from the author upon reasonable request. The analysis was performed on a standard laptop computer with 16 GB RAM and took approximately 2 hours to run all models and generate all figures.

\section{Declaration of AI Usage}

In accordance with AIMS academic integrity policies:

\begin{itemize}
	\item \textbf{Code Generation}: Generative AI (specifically, DeepSeek) was used to assist in writing and debugging R code for data cleaning, analysis, and visualization. All generated code was reviewed, tested, and validated by the author to ensure correctness and reproducibility.
	
	\item \textbf{Writing Assistance}: AI tools were used to assist with grammar checking, phrasing suggestions, and LaTeX formatting. All scientific content, interpretations, and conclusions remain the original work of the author.
	
	\item \textbf{Mathematical Derivations}: All mathematical derivations in this appendix were verified by the author. AI was used to assist with LaTeX formatting of equations, not to generate mathematical content.
	
	\item \textbf{No Plagiarism}: AI was not used to generate text that directly reproduces existing published work without proper citation.
	
	\item \textbf{Supervisor Approval}: The use of AI tools was disclosed to and approved by the supervisor, Dr. Haisa Osmanli.
\end{itemize}

The author takes full responsibility for the accuracy, integrity, and originality of the work presented herein.